\subsection*{2.2.}
现设 $f:X\longrightarrow\Spec(\mathbb F_q)$ 是 $\mathbb F_q$ 上的固有概形.
设 $k$ 是 $\mathbb F_q$ 的代数闭包, 置

\begin{equation}\tag{2.2.1}
\bar X=X\times_{\mathbb F_q}k.
\end{equation}

$\Spec(\mathbb F_q)$ 上的平展层 $R^if_*(\mathbb Z/p\mathbb Z)$
可与视为 $\Gal(k/\mathbb F_q)$ 模的
$H^i_{\mathrm{et}}(\bar X,\mathbb Z/p\mathbb Z)$ 等同. 设

\begin{equation}\tag{2.2.2}
\mathfrak F_q:H^i(X,\mathcal O_X)\longrightarrow H^i(X,\mathcal O_X)
\end{equation}

为下式诱导的线性运算 (它是 $q$ 线性的, 而基域是 $\mathbb F_q$)

\begin{equation}\tag{2.2.3}
\mathfrak F_q:\mathcal O_X\longrightarrow\mathcal O_X,
\qquad \mathfrak F_q(f)=f^q.
\end{equation}

\subsubsection*{命题 2.2.4.}
设 $\varphi_q\in\Gal(k/\mathbb F_q)$ 是“Frobenius”元素,
$\varphi_q(\lambda)=\lambda^q$, 并设 $i\in\mathbb Z$. 有特征多项式恒等式

\begin{equation}\tag{2.2.4.1}
\det_{\mathbb F_p}\!\left(
1-t\varphi_q^{-1}\mid H^i_{\mathrm{et}}(\bar X,\mathbb Z/p\mathbb Z)
\right)
=
\det_{\mathbb F_q}\!\left(
1-t\mathfrak F_q\mid H^i(X,\mathcal O_X)
\right).
\end{equation}

这里左端是 $\pi_1(\Spec(\mathbb F_q))$ 的元素 $\varphi_q^{-1}$
作用于 Galois 模 $H^i_{\mathrm{et}}(\bar X,\mathbb Z/p\mathbb Z)$ 的特征多项式.

命题~2.2.4 是下述结果的推论.

\subsubsection*{命题 2.2.5.}
设 $X$ 在 $\Spec(\mathbb F_q)$ 上固有. 有典范同构

\begin{equation}\tag{2.2.6}
H^i(X,\mathcal O_X)_{\mathrm{ss}}\otimes_{\mathbb F_q}k
\simeq
H^i_{\mathrm{et}}(\bar X,\mathbb Z/p\mathbb Z)
\otimes_{\mathbb F_p}k
\end{equation}

将左端由 $\mathfrak F_q$ 诱导的自同构对应到右端由 $\varphi_q^{-1}$ 诱导的自同构.
后者也可解释为 $(\mathrm{Fr}_q^X)^*\otimes_{\mathbb F_p}k$,
其中 $\mathrm{Fr}_q^X:\bar X\longrightarrow\bar X$ 是 Frobenius $k$ 自同态
(参见 SGA~5, XVIII~2.3.4).

\medskip
\noindent\emph{证明.}
由于

\begin{equation}\tag{2.2.7}
H^i(\bar X,\mathcal O_{\bar X})\simeq
H^i(X,\mathcal O_X)\otimes_{\mathbb F_q}k,
\end{equation}

命题~1.0.10 给出

\begin{equation}\tag{2.2.8}
H^i(\bar X,\mathcal O_{\bar X})_{\mathrm{ss}}\simeq
H^i(X,\mathcal O_X)_{\mathrm{ss}}\otimes_{\mathbb F_q}k.
\end{equation}

因而将 (2.2.8) 与 (2.1.1) 复合即得同构 (2.2.6).

根据 (2.2.7), $q$ 线性运算

\begin{equation}\tag{2.2.9}
\mathfrak F_q:H^i(X,\mathcal O_X)\longrightarrow H^i(X,\mathcal O_X)
\end{equation}

有 $q$ 线性扩张 (参见 1.0.5)

\begin{equation}\tag{2.2.10}
\overline{\mathfrak F}_q:H^i(\bar X,\mathcal O_{\bar X})
\longrightarrow H^i(\bar X,\mathcal O_{\bar X}),
\end{equation}

且也有 $k$ 线性扩张, 仍记为 $\overline{\mathfrak F}_q$,

\begin{equation}\tag{2.2.11}
\overline{\mathfrak F}_q:H^i(\bar X,\mathcal O_{\bar X})
\longrightarrow H^i(\bar X,\mathcal O_{\bar X}).
\end{equation}

设

\begin{equation}\tag{2.2.12}
\varphi_q:H^i(\bar X,\mathcal O_{\bar X})
\longrightarrow H^i(\bar X,\mathcal O_{\bar X})
\end{equation}

为以 $H^i(X,\mathcal O_X)$ 为不动点的 $q$ 线性运算.
在 (2.2.7) 下, 算子 $\mathfrak F_q$、$\varphi_q$ 及 $\overline{\mathfrak F}_q$
也可分别写成 $\mathfrak F_q\otimes\mathrm{id}_k$、$\mathrm{id}\otimes\varphi_q$
及 $\mathfrak F_q\otimes\varphi_q$. 因而它们彼此可交换, 且

\begin{equation}\tag{2.2.13}
\overline{\mathfrak F}_q
=\varphi_q\mathfrak F_q
=\mathfrak F_q\circ\varphi_q.
\end{equation}

它们都与 $\mathfrak F_p$ 可交换, 故 $\mathfrak F_q$ 和 $\varphi_q$ 都作用于
$H^i(\bar X,\mathcal O_{\bar X})^{1-\overline{\mathfrak F}_q}$ 及
$H^i(\bar X,\mathcal O_{\bar X})^{1-\mathfrak F_p}$.
由 (2.2.13), $\mathfrak F_q$ 和 $\varphi_q$ 在
$H^i(\bar X,\mathcal O_{\bar X})^{1-\overline{\mathfrak F}_q}$ 上互为逆.
故 $\mathfrak F_q$ 在

\[
H^i(\bar X,\mathcal O_{\bar X})_{\mathrm{ss}}
\simeq
H^i(\bar X,\mathcal O_{\bar X})^{1-\overline{\mathfrak F}_q}
\otimes_{\mathbb F_q}k
\qquad\text{(参见 1.1.11)}
\]

上的作用与 $\varphi_q\otimes\mathrm{id}_k$ 在

\[
H^i(\bar X,\mathcal O_{\bar X})^{1-\overline{\mathfrak F}_q}
\otimes_{\mathbb F_q}k.
\]

上的作用互为逆. 写作

\begin{equation}\tag{2.2.14}
H^i(\bar X,\mathcal O_{\bar X})^{1-\overline{\mathfrak F}_q}
\simeq
H^i(\bar X,\mathcal O_{\bar X})^{1-\mathfrak F_p}
\otimes_{\mathbb F_p}\mathbb F_q,
\end{equation}

左端的自同构 $\varphi_q$ 通过 (2.2.14) 对应于 $\varphi_q$ 在
$H^i(\bar X,\mathcal O_{\bar X})^{1-\mathfrak F_p}\otimes_{\mathbb F_p}\mathbb F_q$
上诱导的自同构. 因此 $\mathfrak F_q$ 在 $H^i(\bar X,\mathcal O_{\bar X})_{\mathrm{ss}}$
上诱导的自同构, 通过 (2.1.1) 对应于 $\varphi_q^{-1}$ 在
$H^i_{\mathrm{et}}(\bar X,\mathbb Z/p\mathbb Z)\otimes_{\mathbb F_p}k$
上诱导的自同构. 命题~2.2.5 得证.

\medskip
\noindent\emph{底本注.}
证明末尾的印本句子写作 $\varphi_q$, 但命题及紧接其前的互逆关系要求为 $\varphi_q^{-1}$.
英文层遵循数学上相容的读法及校订法文工作稿; 照录法文层保留印本措辞.
